\section{Introduction and statements}

A degeneracy locus remembers where a map loses rank; its scheme
structure can remember substantially more.  In the multiplication
problem studied here, the nilpotent layers of a Fitting scheme recover
geometric data that disappear after reduction.  For Hilbert function
\((1,4,6)\), the first layer recovers a polarized quartic K3 surface.
The next layer detects the first singular Schubert stratum.  The
purpose of this revision is to determine how those two layers fit
together across the rank boundary, including contact collisions and
the first failure of the mixed-regular hypothesis.

Let \(V\) be a complex vector space of dimension
\(e\in\{3,4\}\) and put \(p=e(e-1)/2\).  For an \(e\)-dimensional
relation space
\[
 R\subset\Sym^2V
\]
set
\[
 S_R=\Sym^2V/R,\qquad
 B_R=\C\oplus V\oplus S_R,
\]
where multiplication on \(V\) is the quotient map and
\((V\oplus S_R)S_R=0\).  On
\[
 X_R=\Gr(e,V\oplus S_R)
\]
let \(\mathcal K\) be the tautological subbundle and consider
\begin{equation}\label{eq:failure-intro}
 \widehat D_R=
 V\!\left(
 \Fitt_0\coker[
 \Sym^m(\OO\oplus\mathcal K)\longrightarrow
 B_R\otimes\OO]\right),
 \qquad m\ge2.
\end{equation}
The scheme is independent of \(m\).  Its reduction is the Schubert
divisor \(\widehat\Delta_R\) on which
\(\mathcal K\to V\) drops rank.

Regard \(R\) as a web of quadrics on \(\Pj(V^*)\).  On the nonempty
invariant open \(G_e^\circ\) used below the web is basepoint-free and
the Jacobian hypersurface
\begin{equation}\label{eq:Y-intro}
 Y_R=
 V\!\left(\det(\partial q_i/\partial x_j)\right)
 \subset\Pj(V^*)
\end{equation}
is smooth.  It is a genus-one curve for \(e=3\) and a quartic K3
surface for \(e=4\).

\subsection*{The first nilpotent layer}

Let \(U_R\subset X_R\) be the projection-corank-at-most-one open and
write \(D_R=\widehat D_R\cap U_R\).  The local normal form on this open
is
\[
 \mathcal I_{D_R}=t(t,f)=(t)\cap(t,f)^2,
\]
where \(t=0\) is the Schubert divisor and \(f=0\) is the ramification
equation on it.  If \(\NN_R\) is the nilradical of
\(\OO_{D_R}\), then
\[
 \NN_R^2=0,\qquad
 \Ann(\NN_R)=\mathcal I_{E_R}/\mathcal I_{D_R},
\]
and \(E_R\) is a smooth intrinsic scheme.

\begin{theorem}[Polarized reconstruction]
\label{thm:main-reconstruction}
For \(R\in G_e^\circ\), put
\[
 \LL_R=
 \omega_{E_R}\otimes
 \NN_R^{\otimes-(p+e-1)}.
\]
There is an isomorphism of graded algebras
\[
 \bigoplus_{n\ge0}H^0(E_R,\LL_R^{\otimes n})
 \simeq
 \bigoplus_{n\ge0}H^0(Y_R,\OO_{Y_R}(en)).
\]
Consequently \(D_R\) determines
\((Y_R,\OO_{Y_R}(e))\).  For \(e=4\) it determines
\((Y_R,\OO_{Y_R}(1))\).  The same conclusions hold with
\(\widehat D_R\) in place of \(D_R\); no isomorphism is required to
preserve the ambient Grassmannian, the relation space, or any chosen
projection.
\end{theorem}

\subsection*{The nilpotent filtration across projection rank}

Let \(\widehat\NN_R\) denote the nilradical of
\(\OO_{\widehat D_R}\).  For \(j\ge1\) define
\begin{equation}\label{eq:Zj-intro}
 \ZZ_j(\widehat D_R)
 =
 V\!\left(
 \Ann_{\OO_{\widehat D_R}}
 (\widehat\NN_R^{\,j})\right).
\end{equation}
These schemes are intrinsic and satisfy
\(\ZZ_{j+1}\subset\ZZ_j\).  The next theorem computes the first two
stages together, rather than one stratum at a time.

At a projection-rank-two point let \(H\subset V\) be the image and
\(W=V/H\).  Write
\[
 A_H:\Sym^2H\to S_R
\]
for the pure \(H\)-block and
\[
 \overline B_H:H\otimes W\longrightarrow S_R/\im A_H
\]
for the induced mixed block.  We call \((R,H)\)
\emph{mixed-regular} when \(A_H\) is injective,
\(\overline B_H\) is surjective, and its one-dimensional kernel is
generated by a rank-two tensor.  This condition does not impose
simple contact roots.

On a mixed-regular transverse slice write
\[
 T=\begin{pmatrix}a&b\\c&d\end{pmatrix},\qquad
 \delta=ad-bc,\qquad
 \mathfrak m=(a,b,c,d).
\]
The restricted quartic \(Y_R\cap\Pj(H^\perp)\) is represented by a
binary quartic \(h_H\).  The root-free residual calculation defines a
canonical unsplit contact ideal
\[
 I_H=(\delta,g_0,\ldots,g_4),
\]
whose projectivization on the rank-one Segre quadric is the Cartier
divisor
\[
 V(h_H)\times\Pj^1.
\]

\begin{theorem}[Stratified nilpotent algebra at the first singular
Schubert boundary]
\label{thm:main-depth}
Let \(R\in G_4^\circ\).

\begin{enumerate}
\item On the smooth locus of
\((\widehat D_R)_{\mathrm{red}}\),
\[
 \ZZ_1(\widehat D_R)=E_R,\qquad
 \widehat\NN_R|_{E_R}
 \simeq\OO_{E_R}(-\widehat\Delta_R),
\]
and Theorem~\ref{thm:main-reconstruction} reconstructs the polarized
quartic K3 from this first layer.

\item On every mixed-regular projection-rank-two chart,
\begin{equation}\label{eq:main-contact-factor}
 \mathcal I_{\widehat D_R}
 =
 \delta I_H\cap\mathfrak m^5.
\end{equation}
The factor \(\mathfrak m^5\) is the actual irredundant
\(\mathfrak m\)-primary component throughout the collision boundary.

If
\[
 h_H=\prod_{\nu=1}^s\ell_\nu^{e_\nu},
 \qquad\sum_{\nu=1}^se_\nu=4,
\]
the associated primes are
\[
 (\delta),\ P_1,\ldots,P_s,\ \mathfrak m,
\]
where \(P_\nu\) is the contact prime corresponding to \(\ell_\nu\).
At the generic point of \(P_\nu\), with
\(P_\nu=(\delta,q_\nu)\), the contact-primary type is
\begin{equation}\label{eq:main-collision-primary}
 (\delta^2,\delta q_\nu^{e_\nu},q_\nu^{e_\nu+1}).
\end{equation}
Thus repeated contacts are controlled scheme-theoretically without
choosing four distinct roots.

\item On the same mixed-regular neighbourhood,
\[
 \widehat\NN_R^3=0\ne\widehat\NN_R^2,
\]
and, writing \(\mathcal L\) for the conormal line of the Schubert
divisor and \(\mathcal C_H=V(I_H)\) for the affine contact cone,
\begin{equation}\label{eq:main-grN}
 \gr_{\widehat\NN_R}\OO_{\widehat D_R}
 \simeq
 \OO_{\widehat\Delta_R}
 \oplus
 (\mathcal L\otimes\OO_{\mathcal C_H})
 \oplus
 (\mathcal L^{\otimes2}\otimes\OO_{\Sigma_R}).
\end{equation}
In particular
\[
 \ZZ_1=
 V(I_H\cap\mathfrak m^3),\qquad
 \ZZ_2=\Sigma_R
\]
scheme-theoretically.  On the adjacent projection-corank-one locus,
\(\ZZ_1\) is the ramification support \(E_R\); hence
\eqref{eq:main-grN} describes the specialization of the first layer
to the rank-two boundary and identifies the second layer there.

\item If \(h_H\equiv0\), equivalently the pencil
\(\Pj(H^\perp)\) is a line contained in \(Y_R\), then
\begin{equation}\label{eq:main-containment}
 I_H=(\delta),\qquad
 \mathcal I_{\widehat D_R}
 =\delta^2\mathfrak m
 =(\delta^2)\cap\mathfrak m^5.
\end{equation}
The containment locus in \(\Gr(2,V)\) is canonically
\(F_1(Y_R)\), the finite reduced Fano scheme of lines on the smooth
quartic.

\item Suppose \(A_H\) is injective and \(\overline B_H\) remains
surjective with one-dimensional kernel, but that kernel becomes a
decomposable tensor.  Under the transverse nonvanishing condition of
Definition~\ref{def:decomposable-kernel-regular}, the nilpotent index
is still three while
\[
 \ZZ_2=
 V(c,d,a^2,ab,b^2).
\]
Thus the second intrinsic depth support becomes a length-three
nonreduced thickening of the rank-two vertex on the first
non-mixed-regular tensor-rank wall.
\end{enumerate}

At every projection rank, independently of these hypotheses, the full
Fitting ideal has the exact block presentation of
Theorem~\ref{thm:all-corank-presentation}.  The presentation theorem
and the structure statements above are logically distinct.
\end{theorem}

The theorem is proved in
Sections~\ref{sec:depth}--\ref{sec:stratified-nilpotent}.  Its central
identity is
\[
 J=I_H\cap\mathfrak m^3,\qquad
 \delta J=\delta I_H\cap\mathfrak m^5.
\]
The first equality identifies the contact saturation hidden in the
root-free residual minors; the second converts it into a genuine
primary statement at the rank-two vertex.  It also explains why a
collision alters the contact-primary type but leaves the vertex
component fixed.  The associated-graded formula
\eqref{eq:main-grN} then follows by cancelling successive powers of
\(\delta\).

The result separates three levels of generality.  The residual block
presentation holds at every projection rank.  The complete contact
specialization and associated-graded algebra hold on the entire
mixed-regular rank-two locus, including repeated roots and the
containment locus.  The first tensor-rank boundary beyond
mixed-regularity is described by the decomposable-kernel theorem.
This separation is used throughout the paper.

\subsection*{The contact incidence}

Put \(G=\Gr(2,V)\), with universal quotient bundle \(\mathcal Q\).
The quartic equation of \(Y_R\) restricts to a global section
\[
 h_R\in H^0(G,\Sym^4\mathcal Q).
\]
The incidence
\[
 \RR_R=
 \{(H,[\alpha])\in G\times Y_R:
                 H\subset\ker\alpha\}
\]
is a \(\Pj^2\)-bundle over \(Y_R\), hence smooth and irreducible.
Over the locus where \(h_H\) is not identically zero it is finite
flat of degree four over \(G\).  Its global discriminant is
\[
 \Disc(h_R)\in H^0(G,(\det\mathcal Q)^{12}),
\]
so the discriminant divisor has Pluecker class \(12H\).  The
interpretation of the four etale sheets as the four simple-contact
primary branches requires, in addition, the mixed-regular rank
conditions.  On the containment locus the fibre is a projective line
and the finite-cover interpretation is replaced by
Theorem~\ref{prop:containment-lines}.

\subsection*{Simple contacts as the split form}

On the simple-contact mixed-regular open the unsplit theorem recovers
the earlier etale primary decomposition.

\begin{theorem}[Corank-two primary structure on the simple-contact
open]\label{thm:main-coranktwo}
There is a nonempty invariant open
\(G_4^\dagger\subset G_4^\circ\) such that, for each
\(R\in G_4^\dagger\), a dense open set of rank-two images
\(H\subset V\) has the following property.  Near every plane with
image \(H\), after an etale extension and smooth auxiliary
coordinates,
\[
 \mathcal I_{\widehat D_R}
 =
 (\delta)\cap P_1^2\cap P_2^2\cap P_3^2\cap P_4^2
 \cap\mathfrak m^5.
\]
The rank-two stratum is an additional associated support,
\[
 \widehat\NN_R^3=0\ne\widehat\NN_R^2,\qquad
 \Ann(\widehat\NN_R^2)
 =\mathfrak m/\mathcal I_{\widehat D_R}.
\]
\end{theorem}

Theorem~\ref{thm:main-depth} is the unsplit and stratified extension:
simple roots give \(e_\nu=1\) in
\eqref{eq:main-collision-primary}; collisions give higher \(e_\nu\);
and \(h_H\equiv0\) gives the separate containment algebra
\eqref{eq:main-containment}.

\subsection*{Variation and finite ambiguity}

The polarized K3s vary nontrivially.

\begin{theorem}[Variation and finite ambiguity]
\label{thm:main-moduli}
For \(e=4\), the image of
\[
 R\longmapsto(Y_R,\OO_{Y_R}(1))
\]
in quartic K3 moduli has dimension nine.  On a nonempty invariant
open, the induced map from projective equivalence classes of relation
spaces to its image is generically finite.  Thus, over a general
point of this image, \(D_R\) and \(\widehat D_R\) determine at most
finitely many isomorphism classes of algebras \(B_R\).
\end{theorem}

The exact rank-\(24\) differential witness in
Section~\ref{sec:moduli} proves this assertion for the actual Jacobian
quartic map.  Section~\ref{sec:reye-priority} separates that arithmetic
witness from the classical geometric reason for the dimension count.
A regular web gives the classical Reye Enriques surface and its K3
cover.  The web parameter is \(24\)-dimensional and the infinitesimal
projective orbit is \(15\)-dimensional, leaving the classical
nine-dimensional quotient.  We make the parameter-to-Hilbert bridge
explicit and isolate the finite involution ambiguity on a fixed
polarized K3.

The number nine is therefore not used as a novelty signal.  The
intrinsic statement is that an abstract nonreduced multiplication
scheme recovers the polarized K3 and, in addition, the stratified
nilpotent algebra described in Theorem~\ref{thm:main-depth}.

\subsection*{Classical and priority boundaries}

The bilinear incidence K3, the fixed-point-free factor-exchange
involution, and the associated Reye congruence are classical.  We
compare them with Cossec, Arrondo--Sols, Dolgachev--Reider,
Ingalls--Kuznetsov, and the modern nodal-Enriques treatment in
Section~\ref{sec:reye-priority}.  We likewise compare the nilpotent
filtration with ribbons and multiple structures.  The claim is not
that nilpotent filtrations in general are new; it is that the
particular Fitting algebra in \eqref{eq:failure-intro} has intrinsic
graded pieces that recover the polarized K3, the unsplit contact
Cartier geometry, and the first singular Schubert support.

Ballico's failure-locus work is treated as a genuine predecessor.
The accessible 1996 continuation is compared at theorem level.  The
1993 article is cited explicitly and is not used as a premise in any
non-anticipation argument.  The mathematical theorems of this paper
are stated and proved independently of that documentary question.
